\section{Discussion}
\label{sec:discussion}

\paragraph{The contribution is a measurement protocol, not a converter leaderboard.}
The main result is not that noMDL conversion is always harmless, faster, or
better than the NVIDIA Asset Converter baseline. The contribution is to make
material conversion visible to VLM and embodied-language evaluation: proxy
metrics, VLM answers, point grounding, coordinate compliance, material-effect
risk, and embodied-stack usability answer different reliability questions.

\paragraph{Grounding deltas are only interpretable after the prompt contract is fixed.}
The GRScenes stress result shows why proxy-similar renders are not enough:
Gemma4 follows the normalized-1000 contract, while Qwen2.5-VL exposes stronger
raw-pixel behavior under the same prompt family. The finding is about
measurement contracts, not model ranking; coordinate frames and parser coverage
must be fixed before a material-conversion delta is interpretable.

\paragraph{Material and embodied evidence need explicit boundaries.}
The material-effect audit prevents a common failure mode in conversion papers:
counting static conversion success, zero-MDL outputs, or effect-token coverage
as visual or task success. The four GRScenes-covered bins support bounded
selected qualitative evidence; clearcoat is a selected NVIDIA failure case, and
procedural texture is a limitation for both converted conditions. Likewise, the
99-episode official KuJiaLe InternNav run is sanity evidence over three local
official scenes, not broad embodied robustness evidence.
